\begin{center}
    \textbf{Boa Esperança}
\end{center}

\noindent
Por mais que você corra irmão\\
Pra sua guerra vão nem se lixar\\
Esse é o xis da questão\\
Já viu eles chorar pela cor do orixá?\\
E os camburão o que são?\\
Negreiros a retraficar\\
Favela ainda é senzala jão\\
Bomba relógio prestes a estourar\\
(x2)

\vspace{0.3cm}

\noindent
Aí\\
O tempero do mar foi lágrima de preto\\
Papo reto, como esqueletos, de outro dialeto\\
Só desafeto, vida de inseto, imundo\\
Indenização? Fama de vagabundo\\
Nação sem teto, Angola, keto, congo, soweto\\
A cor de Eto'o, maioria nos gueto\\
Monstro sequestro, capta-três, rapta\\
Violência se adapta, um dia ela volta pu cêis\\
Tipo campos de concentração, prantos em vão\\
Quis vida digna, estigma, indignação\\
O trabalho liberta, ou não\\
Com essa frase quase que os nazi, varre os judeu? extinção\\
[...]


\noindent
Aí\\
Nessa equação, chata, polícia mata -- Plow!\\
Médico salva? Não! Por que? Cor de ladrão\\
Desacato invenção, maldosa intenção\\
Cabulosa inversão, jornal distorção\\
Meu sangue na mão dos radical cristão\\
Transcendental questão, não choca opinião\\
Silêncio e cara no chão, conhece?\\
Perseguição se esquece? Tanta agressão enlouquece\\
Vence o Datena, com luto e audiência\\
Cura baixa escolaridade com auto de resistência\\
Pois na era cyber, cêis vai ler\\
Os livro que roubou nosso passado igual alzheimer, e vai ver\\
Que eu faço igual burkina faso\\
Nóiz quer ser dono do circo\\
Cansamos da vida de palhaço\\
É tipo moisés e os hebreus, pés no breu\\
Onde o inimigo é quem decide quando ofendeu\\
(Cê é loco meu)\\
[...]

\vspace{0.3cm}

\begin{center}
\footnotesize
(Disponível em: \url{https://www.vagalume.com.br/emicida/boa-esperanca.html}.
Acesso em: 08 out. 2017. Adaptado.)
\end{center}

\textbf{COLUNI 2018 -- Questão 04}\par\medskip O rap Boa esperança e o livro Quarto de Despejo de Carolina Maria de Jesus, apesar de terem sido escritos em épocas diferentes, dialogam entre si no que se refere ao modo de ver a favela, porque a definem como:\begin{enumerate}
\item um lugar no qual a fome e a miséria governam a vida das pessoas.
\item um espaço que leva à degradação da fé de seus moradores.
\item um espaço que abriga os excluídos e explorados pela sociedade.
\item um lugar onde vivem, principalmente, pessoas desempregadas.\end{enumerate}
